\section{生态职责：平台、技术提供者、报告机构与调查人员}

最后回到端到端系统。儿童安全生态中没有一个参与者拥有全部能力：平台掌握用户关系和产品 action，技术提供者提供 detection/review primitives，法定报告机构接收并丰富线索，执法与受害者服务机构依法调查和保护。责任清晰不是推卸，而是避免模型供应商越权、平台外包治理或调查人员被低质量数据淹没。

\subsection{责任边界与反馈闭环}

平台应拥有 policy、用户通知、申诉、账户处置与法定义务；Thorn 等技术提供者应提供经验证、可审计的工具并保护敏感数据；NCMEC 等报告机构按法定角色处理报告和线索；执法负责调查、证据与法律程序；受害者支持机构提供持续保护。Verified outcome 只有在授权、最小化和目的限制下才能反馈给模型与 reference system。

\lecturefigure{16-ecosystem-roles.png}{儿童安全结果依赖平台、技术提供者、报告机构与调查人员的清晰分工。}{本地字幕 37:10--37:53；基于课堂 closing synthesis 的概念重绘。}

\begin{knowledgebox}{读图：每次 handoff 都需要 contract}
Platform 到 provider 的 contract 规定输入、输出、SLA、数据用途和 incident；platform 到 reporting agency 的 contract 规定法定字段、重复处理和证据链；agency 到 investigators 的 handoff 需要 prioritization 与 jurisdiction；investigation outcome 回流时必须限制用途。模糊 handoff 会产生重复报告、字段缺失、无人处理或越权访问。
\end{knowledgebox}

\begin{importantbox}{Mission-driven business model 的工程含义}
非营利使命与商业软件并不天然冲突：收费产品可为持续研发、支持、审计和高可用基础设施提供资金；使命治理则约束客户选择、数据用途和产品优先级。真正的验收标准不是组织标签，而是能否长期维护安全工具、公开边界、接受监督，并让保护结果而非 engagement 成为北极星指标。
\end{importantbox}

\teachervoice{Cordua 的收尾把角色说得很直接：平台要检测和采取行动，调查人员要识别受害者、追查案件。技术的作用是把双方之间的信息变得更快、更少重复、更可行动，而不是让一个模型替代制度与专业人员。}

\subsection{本章小结}

儿童安全是跨组织系统。清晰 role、data contract、audit trail 和反馈边界能让各方发挥专长；模糊责任则会把风险转移给最缺乏上下文或最缺乏资源的一环。

